\chapter{Advanced Features of \textsc{Lark}}

\begin{figure}[h] 
  \centering
    \includegraphics[width=17cm]{Abbildungen/larkhall.jpg}
  \caption{Larkhall in South Lanarkshire, Scotland.}
  \label{fig:larkhall.jpg}
  \end{figure}

When we discussed the parser generator \simtextsc{Lark} in chapter 7, we were not able to discuss all the features of
\simtextsc{Lark}, since the underlying theory, which is the theory of \simtextsc{Lalr} parsing, had not been developed.
Since we have acquainted ourselves with this theory in the previous chapter, we are finally able to discuss some of the
more advanced features of \simtextsc{Lark}.
\begin{enumerate}
\item First, we show how shift-reduce and reduce-reduce conflicts are reported by \simtextsc{Lark} and how the
      \simtextsc{Lalr} states that \simtextsc{Lark} computes can be inspected.
\item Then we discuss how these conflicts can be resolved.  As \simtextsc{Lark} resolves every shift-reduce
      conflict in favour of shifting and offers no way to change this decision, the precedences and
      associativities of the operators have to be encoded into the grammar itself.
\end{enumerate}
All examples of this chapter are available as Jupyter notebooks in the following directory:
\\[0.2cm]
\hspace*{1.3cm}
\href{https://github.com/karlstroetmann/Formal-Languages/tree/master/Python/Chapter-10-Lark}{\texttt{Python/Chapter-10-Lark}}.
\vspace*{-0.2cm}

\begin{center}
  \begin{tabular}[t]{|l|l|}
    \hline
    notebook & topic \\
    \hline
    \hline
    \href{https://github.com/karlstroetmann/Formal-Languages/tree/master/Python/Chapter-10-Lark/01-Conflicts.ipynb}{\texttt{01-Conflicts.ipynb}}            & shift-reduce conflicts and how to inspect the states \\
    \hline
    \href{https://github.com/karlstroetmann/Formal-Languages/tree/master/Python/Chapter-10-Lark/02-Conflicts-Resolved.ipynb}{\texttt{02-Conflicts-Resolved.ipynb}}   & encoding operator precedences into the grammar       \\
    \hline
    \href{https://github.com/karlstroetmann/Formal-Languages/tree/master/Python/Chapter-10-Lark/03-Look-Ahead.ipynb}{\texttt{03-Look-Ahead.ipynb}}           & a reduce-reduce conflict caused by insufficient look-ahead \\
    \hline
    \href{https://github.com/karlstroetmann/Formal-Languages/tree/master/Python/Chapter-10-Lark/04-Look-Ahead-Solved.ipynb}{\texttt{04-Look-Ahead-Solved.ipynb}}    & the same grammar, rewritten so that it is \simtextsc{Lalr}(1) \\
    \hline
    \href{https://github.com/karlstroetmann/Formal-Languages/tree/master/Python/Chapter-10-Lark/05-Mysterious-Conflicts.ipynb}{\texttt{05-Mysterious-Conflicts.ipynb}} & a mysterious reduce-reduce conflict                  \\
    \hline
  \end{tabular}
\end{center}

\section{Shift-Reduce and Reduce-Reduce Conflicts}
In this section we show how shift-reduce and reduce-reduce conflicts are dealt with in \simtextsc{Lark}.
Figure \ref{fig:Conflicts.ipynb} on page \pageref{fig:Conflicts.ipynb} shows a grammar for arithmetical
expressions that is ambiguous because it does not specify the precedence of the different arithmetical
operators.  

\begin{figure}[!ht]
\centering
\begin{minted}[ frame        = lines, 
                framesep     = 0.3cm, 
                bgcolor      = sepia,
                numbers      = left,
                numbersep    = -0.2cm,
                xleftmargin  = 0.8cm,
                xrightmargin = 0.8cm,
              ]{text}
    expr : expr "+" expr  -> add
         | expr "*" expr  -> mul
         | NUMBER         -> number

    NUMBER : /0|[1-9][0-9]*/

    %import common.WS
    %ignore WS
\end{minted}
\vspace*{-0.3cm}
\caption{An ambiguous grammar for arithmetical expressions.}
\label{fig:Conflicts.ipynb}
\end{figure}
\FloatBarrier

\noindent
Note that in a \simtextsc{Lark} grammar the syntactical variables are written in lower case, while the names of
the token classes are written in upper case.  The operators \squoted{+} and \squoted{*} occur as
\blue{anonymous terminals}: they are given as string literals inside the grammar rules and \simtextsc{Lark}
invents the names \texttt{PLUS} and \texttt{STAR} for them.  These invented names show up in the messages that
we discuss below.  The three alternatives additionally carry the \blue{aliases}
\texttt{add}, \texttt{mul} and \texttt{number}, which are attached with the operator ``\texttt{->}''.  These
aliases only determine the names of the nodes of the parse tree, they have no influence on the parse table.
They are discussed in the remark on page \pageref{remark:tree-shape}. 

This grammar does not specify whether the string
\\[0.2cm]
\hspace*{1.3cm} 
``\texttt{1 + 2 * 3}'' \quad is interpreted as \quad  ``\texttt{(1 + 2) * 3}'' \quad or as \quad ``\texttt{1 +
  (2 * 3)}''. 
\\[0.2cm]
Since every \simtextsc{Lalr} grammar is unambiguous, but the grammar shown in Figure \ref{fig:Conflicts.ipynb} is
ambiguous, it has to have at least one shift-reduce or reduce-reduce conflict.  This grammar is part of the
jupyter notebook  
\\[0.2cm]
\hspace*{1.3cm}
\href{https://github.com/karlstroetmann/Formal-Languages/tree/master/Python/Chapter-10-Lark/01-Conflicts.ipynb}{Python/Chapter-10-Lark/01-Conflicts.ipynb}.
\\[0.2cm]

\subsection{Making \textsc{Lark} Talk About Its Conflicts}
By default, \simtextsc{Lark} is remarkably tight-lipped about shift-reduce conflicts: the parser is built, the
conflicts are resolved silently, and nothing is printed.  There are two reasons for this silence.
\begin{enumerate}
\item The conflict messages are written to the logger object \texttt{lark.logger}, and the level of this logger
      is set to \texttt{logging.CRITICAL} when \simtextsc{Lark} is imported.  Therefor only critical errors
      are printed but warnings are discarded.
\item The messages are only issued with the severity \texttt{warning} if the keyword argument
      \texttt{debug=True} is given when the parser is constructed.  Otherwise they are issued with the severity
      \texttt{debug}.
\end{enumerate}
Therefore, in order to see the conflicts of the grammar given in Figure \ref{fig:Conflicts.ipynb}, we have to
create the parser as shown in Figure \ref{fig:conflicts-code} on page \pageref{fig:conflicts-code}.

\begin{figure}[!ht]
\centering
\begin{minted}[ frame        = lines, 
                framesep     = 0.3cm, 
                bgcolor      = sepia,
                numbers      = left,
                numbersep    = -0.2cm,
                xleftmargin  = 0.3cm,
                xrightmargin = 0.3cm,
              ]{python3}
    import logging
    from lark import Lark, logger

    logger.setLevel(logging.DEBUG)

    parser = Lark(grammar, start='expr', parser='lalr', debug=True)
\end{minted}
\vspace*{-0.3cm}
\caption{Creating a parser so that conflicts are reported.}
\label{fig:conflicts-code}
\end{figure}
\FloatBarrier

\noindent
When we run this code with the grammar of Figure \ref{fig:Conflicts.ipynb}, we get the output that is shown in
Figure \ref{fig:Conflicts.warnings} on page \pageref{fig:Conflicts.warnings}.

\begin{figure}[!ht]
\centering
\begin{minted}[ frame        = lines, 
                framesep     = 0.3cm, 
                bgcolor      = sepia,
                numbers      = left,
                numbersep    = -0.2cm,
                xleftmargin  = 0.3cm,
                xrightmargin = 0.3cm,
              ]{text}
    Shift/Reduce conflict for terminal PLUS: (resolving as shift)
     * <expr : expr PLUS expr>
    Shift/Reduce conflict for terminal STAR: (resolving as shift)
     * <expr : expr PLUS expr>
    Shift/Reduce conflict for terminal PLUS: (resolving as shift)
     * <expr : expr STAR expr>
    Shift/Reduce conflict for terminal STAR: (resolving as shift)
     * <expr : expr STAR expr>
\end{minted}
\vspace*{-0.3cm}
\caption{The conflicts reported for the grammar of Figure \ref{fig:Conflicts.ipynb}.}
\label{fig:Conflicts.warnings}
\end{figure}
\FloatBarrier

\noindent
Hence there are four shift-reduce conflicts.  Every one of these messages names
\begin{enumerate}
\item the terminal $o$ for which the entry $\textsl{action}(s,o)$ of the action table is ambiguous, and
\item the grammar rule $r$ that could have been used to reduce the symbol stack in this situation.
\end{enumerate}
Unfortunately, the message does not name the state $s$ in which the conflict occurs, and it does not show the
follow sets of the rules involved.  We will remedy this in the next subsection.

The important observation is the parenthesised remark ``\blue{\texttt{(resolving as shift)}}'': \simtextsc{Lark} resolves
\emph{every} shift-reduce conflict in favour of shifting.  There is no way to change this decision for an
individual conflict.  The only choice we have is to declare that we do not want to tolerate conflicts at all.
If the parser is created with the keyword argument \blue{\texttt{strict=True}}, then the first shift-reduce conflict
raises an exception of class \texttt{GrammarError} whose message is
\\[0.2cm]
\hspace*{1.3cm}
\texttt{Shift/Reduce conflict for terminal PLUS. [strict-mode]}
\\
\hspace*{1.3cm}
\texttt{ * <expr : expr PLUS expr>}
\\[0.2cm]
It is a good idea to develop a grammar with \texttt{strict=True}: this way, a conflict cannot be overlooked.
Note that \texttt{strict=True} additionally checks the terminals for collisions and that this check needs the
package \texttt{interegular}, which is not installed together with \simtextsc{Lark} by default.  If it is
missing, an exception of class \texttt{LexError} is raised as soon as the grammar is free of conflicts.

Reduce-reduce conflicts are treated much more severely.  A reduce-reduce conflict always raises an exception of
class \texttt{GrammarError}, regardless of the value of \texttt{strict}.  We will see an example of this in
Section \ref{section:resolving-conflicts}.

\subsection{Inspecting the \textsc{L\small{ALR}} States}
\simtextsc{Lark} provides no method that would print the \simtextsc{Lalr} states it has computed.  However, if
the parser is created with \texttt{debug=True}, then the parse table is kept in a form in which the states are
still the sets of marked rules that we have discussed in the previous chapter, so we can print them ourselves.
Figure \ref{fig:dump-states} on page \pageref{fig:dump-states} shows a function that does this.  The same
function is used in all notebooks of this chapter.

\begin{figure}[!ht]
\centering
\begin{minted}[ frame        = lines, 
                framesep     = 0.3cm, 
                bgcolor      = sepia,
                numbers      = left,
                numbersep    = 0.4cm,
                xleftmargin  = 0.0cm,
                xrightmargin = 0.0cm,
              ]{python3}
from lark.parsers.lalr_analysis import Shift

def dump_states(lark_parser):
  table  = lark_parser.parser.parser.parser.parse_table
  key    = lambda state: sorted(str(rule) for rule in state)
  start  = next(iter(table.start_states.values()))
  rest   = sorted((s for s in table.states if s != start), key=key)
  order  = [start] + rest
  number = { state: i for i, state in enumerate(order) }
  for state in order:
    print(f'state {number[state]}')
    for marked_rule in sorted(state, key=str):
      print(f'    {marked_rule}')
    print()
    for token, (action, arg) in sorted(table.states[state].items()):
      if action is Shift:
        print(f'    {token:<8} shift and go to state {number[arg]}')
      else:
        print(f'    {token:<8} reduce using rule {arg}')
    print()
\end{minted}
\vspace*{-0.3cm}
\caption{A function that prints the \simtextsc{Lalr} states computed by \simtextsc{Lark}.}
\label{fig:dump-states}
\end{figure}
\FloatBarrier

\noindent
The chain of attributes \texttt{lark\_parser.parser.parser.parser} looks confusing, but it is easily explained: the
object returned by \texttt{Lark} delegates to a \blue{parsing frontend}, which delegates to an object of class
\texttt{LALR\_Parser}, which finally contains the object of class \texttt{\_Parser} that stores the attribute
\texttt{parse\_table}.  The states of this table are \emph{frozen sets of marked rules}, which is exactly the
notion that we have developed in the previous chapter.  They are not numbered, so the dictionary
\texttt{number} assigns a number to every state in order to make the output readable.  Since the states are
sets, iterating over them is not reproducible from one run to the next.  Therefore, the states are sorted
before they are numbered and the start state is put first, so that it receives the number $0$.

Given the grammar shown in Figure \ref{fig:Conflicts.ipynb}, \simtextsc{Lark} creates seven different states.
Figure \ref{fig:Conflicts.ipynb:state5} on page \pageref{fig:Conflicts.ipynb:state5} shows the most interesting
of these, together with its actions.

\begin{figure}[!ht]
\centering
\begin{minted}[ frame        = lines, 
                framesep     = 0.3cm, 
                bgcolor      = sepia,
                numbers      = left,
                numbersep    = -0.2cm,
                xleftmargin  = 0.8cm,
                xrightmargin = 0.8cm,
              ]{text}
    state 5
        <expr : expr * PLUS expr>
        <expr : expr * STAR expr>
        <expr : expr PLUS expr * >

        $END     reduce using rule <expr : expr PLUS expr>
        PLUS     shift and go to state 2
        STAR     shift and go to state 3
\end{minted}
\vspace*{-0.3cm}
\caption{One of the states with a conflict.}
\label{fig:Conflicts.ipynb:state5}
\end{figure} %$


\noindent
Note that \simtextsc{Lark} marks the position of the parser inside a rule with the character
``\texttt{*}'' instead of the bullet ``$\bullet$'' that we have been using so far.  The state shown in Figure
\ref{fig:Conflicts.ipynb:state5} is reached after an expression of the form
$\textsl{expr}\; \aquoted{+} \;\textsl{expr}$ has been read.  There are two shift-reduce conflicts in this
state: both for the token \texttt{PLUS} and for the token \texttt{STAR}, the parser could either reduce with
the rule
\\[0.2cm]
\hspace*{1.3cm}
$\textsl{expr} \rightarrow \textsl{expr} \aquoted{+} \textsl{expr}$
\\[0.2cm]
or shift the operator.  We see that both conflicts have been decided in favour of shifting: the only reduce
action that is left in this state is the one for the end of input, which is called \texttt{\$END} by
\simtextsc{Lark}.  Note that the printed state no longer shows the discarded reduce actions, nor does it show
the follow sets of the marked rules: the parse table only records the decision that has actually been taken.
If we want to know which actions have been discarded, we have to read the warnings shown in Figure
\ref{fig:Conflicts.warnings}.

Of course, resolving these conflicts in favour of shifting is wrong: the string
\\[0.2cm]
\hspace*{1.3cm}
\texttt{1 * 2 + 3}
\\[0.2cm]
is then interpreted as ``\texttt{1 * (2 + 3)}'' and not as ``\texttt{(1 * 2) + 3}''.  You can convince yourself
of this by inspecting the parse tree that \simtextsc{Lark} produces for this string.

\section{Operator Precedence \label{section:operator-precedence}}

\begin{figure}[!ht]
  \begin{center}    
  \framebox{
  \framebox{
    \begin{minipage}[t]{9cm}    
  \begin{eqnarray*}
  \textsl{expr}    & \rightarrow & \;\textsl{expr} \aquoted{+} \textsl{expr}  \\
                   & \mid        & \;\textsl{expr} \aquoted{-} \textsl{expr}  \\
                   & \mid        & \;\textsl{expr} \aquoted{*} \textsl{expr}  \\
                   & \mid        & \;\textsl{expr} \aquoted{/} \textsl{expr}  \\
                   & \mid        & \;\textsl{expr} \aquoted{\symbol{94}} \textsl{expr}  \\
                   & \mid        & \aquoted{(} \textsl{expr} \aquoted{)}       \\
                   & \mid        & \;\simtextsc{Number}                            
  \end{eqnarray*}
  \vspace*{-0.5cm}
  \end{minipage}}}
  \end{center}
  \caption{An ambiguous grammar for arithmetical expressions.}
  \label{fig:grammar-resolved.g}
\end{figure}

\noindent
Consider the grammar for arithmetical expressions shown in Figure \ref{fig:grammar-resolved.g} on page
\pageref{fig:grammar-resolved.g}.  This grammar is ambiguous for the same reason as the grammar discussed in
the previous section: it does not specify the precedences and associativities of the operators.

\begin{tcolorbox}[
  enhanced,                      % Enables advanced features like borderlines
  width=0.95\textwidth,          % Sets the width 
  center,                        % Centers the box on the page
  colback=yellow!20,             % A light yellow background for the text
  colframe=blue,                 % Blue color for the main outer frame
  boxrule=1pt,                   % Thickness of the outer frame
  borderline={1pt}{3pt}{blue},   % Draws the inner second frame (1pt thick, 3pt inward)
  sharp corners,                 % Keeps the square corners (like \framebox)
  top=0.5cm,                     % 0.5 cm distance from text to top frame
  bottom=0.5cm,                  % 0.5 cm distance from text to bottom frame
  left=0.5cm,                    % 0.5 cm distance from text to left frame
  right=0.5cm,                   % 0.5 cm distance from text to right frame
  boxsep=0pt                     % Resets default internal spacing so your 0.5cm is exact
]
  \textbf{There is no way to declare the precedence of an operator in \simtextsc{Lark}.}
  Neither a precedence nor an associativity can be attached to a terminal, and the resolution of an
  individual shift-reduce conflict cannot be influenced. Every shift-reduce conflict is resolved in favour
  of shifting.
\end{tcolorbox}

\noindent
As every shift-reduce conflict is resolved in favour of shifting, all operators effectively have the same
precedence and, moreover, they all associate to the right, since shifting an operator of the same precedence
is exactly what a right associative operator requires.  This explains the parse trees that we have seen in the
previous section: the string ``\texttt{1*2+3}'' is read as ``\texttt{1*(2+3)}'' because an operator
that groups to the right collects everything that follows it.

Hence, if we want the operators to have different precedences, we have to encode these precedences into the
grammar itself.  This is done by introducing one syntactical variable for every precedence level.
The technique is the one that we have already used in the chapter on context-free languages, and it should by
now be familiar: the variable of a given precedence level is defined in terms of the variable of the next
higher precedence level.
\begin{itemize}
\item If an operator of precedence level $n$ is \blue{left associative}, the corresponding rule is
      \blue{left recursive}:
      \\[0.2cm]
      \hspace*{1.3cm}
      $v_n \rightarrow v_n \;o\; v_{n+1}$.
\item If the operator is \blue{right associative}, the rule is \blue{right recursive}:
      \\[0.2cm]
      \hspace*{1.3cm}
      $v_n \rightarrow v_{n+1} \;o\; v_n$.
\item If the operator is \blue{non-associative}, the rule uses neither recursion:
      \\[0.2cm]
      \hspace*{1.3cm}
      $v_n \rightarrow v_{n+1} \;o\; v_{n+1}$.
\end{itemize}
Figure \ref{fig:grammar-layered} on page \pageref{fig:grammar-layered} shows the grammar of Figure
\ref{fig:grammar-resolved.g} rewritten in this way.  The variable \texttt{expr} corresponds to precedence level
1, \texttt{prod} to precedence level 2, and \texttt{fact} to precedence level 3.  The variable \texttt{atom}
collects the expressions that bind most tightly, that is, numbers and parenthesised expressions.

\begin{figure}[!ht]
\centering
\begin{minted}[ frame        = lines, 
                framesep     = 0.3cm, 
                bgcolor      = sepia,
                numbers      = left,
                numbersep    = -0.2cm,
                xleftmargin  = 0.8cm,
                xrightmargin = 0.8cm,
              ]{text}
    ?expr : expr "+" prod  -> add
          | expr "-" prod  -> sub
          | prod

    ?prod : prod "*" fact  -> mul
          | prod "/" fact  -> div
          | fact

    ?fact : atom "^" fact  -> power
          | atom

    ?atom : "(" expr ")"
          | NUMBER         -> number

    NUMBER : /0|[1-9][0-9]*/

    %import common.WS
    %ignore WS
\end{minted}
\vspace*{-0.3cm}
\caption{An unambiguous \simtextsc{Lark} grammar for arithmetical expressions.}
\label{fig:grammar-layered}
\end{figure}
\FloatBarrier

\noindent
Observe how the associativities are expressed: the rules for \texttt{expr} and \texttt{prod} are left recursive,
because the additive and multiplicative operators associate to the left, while the rule
\\[0.2cm]
\hspace*{1.3cm}
\texttt{fact : atom "\symbol{94}" fact}
\\[0.2cm]
is right recursive, because exponentiation associates to the right.  This grammar is an \simtextsc{Lalr}(1)
grammar: if we build a parser for it with \texttt{strict=True}, no exception is raised and no warning is
printed.  The Jupyter notebook
\\[0.2cm]
\hspace*{1.3cm}
\href{https://github.com/karlstroetmann/Formal-Languages/tree/master/Python/Chapter-10-Lark/02-Conflicts-Resolved.ipynb}{Python/Chapter-10-Lark/02-Conflicts-Resolved.ipynb}
\\[0.2cm]
contains this grammar.  Checking a few examples confirms that the generated parser behaves as intended:
\begin{enumerate}[(a)]
\item ``\texttt{1-2-3}'' is parsed as ``\texttt{(1-2)-3}'', 
\item ``\texttt{1-2*3}'' is parsed as ``\texttt{1-(2*3)}'', and 
\item ``\texttt{2\symbol{94}3\symbol{94}4}'' is parsed as ``\texttt{2\symbol{94}(3\symbol{94}4)}''.  
\end{enumerate}


\remarkEng\label{remark:tree-shape}
Encoding the precedences into the grammar has a price: the parse trees get deeper, since every expression is
now wrapped into one node per precedence level.  This is a nuisance if the parse tree is processed further.
\simtextsc{Lark} offers two remedies for this.
\begin{enumerate}
\item If the name of a syntactical variable is prefixed with a question mark, as in
      \texttt{?expr}, then a node of this variable is removed from the tree whenever it has exactly one child.
      This inlines the chains \texttt{expr} $\rightarrow$ \texttt{prod} $\rightarrow$ \texttt{fact}
      $\rightarrow$ \texttt{atom} that carry no information.
\item Individual alternatives can be given a name using the operator ``\texttt{->}'', as in
      \\[0.2cm]
      \hspace*{1.3cm}
      \texttt{?expr : expr "+" prod -> add}
      \\[0.2cm]
      so that the resulting tree node is labelled \texttt{add} instead of \texttt{expr}.
\end{enumerate}
Note that an alternative carrying an alias is never inlined, even if it has a single child.  This is why the
alternative \texttt{NUMBER -> number} survives although \texttt{?atom} begins with a question mark.

The notebook uses both devices: it writes \texttt{?expr}, \texttt{?prod}, \texttt{?fact} and \texttt{?atom},
and it labels the alternatives with the aliases \texttt{add}, \texttt{sub}, \texttt{mul}, \texttt{div},
\texttt{power} and \texttt{number}.  An object of class \texttt{Transformer} that provides one method per
alias then converts the parse tree into the nested tuples that we have been using as abstract syntax trees.

Both features concern the shape of the parse tree only, they have no influence on the parse table and
therefore cannot introduce or remove conflicts.
\eox

\section{Resolving Shift-Reduce and Reduce-Reduce Conflicts \label{section:resolving-conflicts}}
We start our discussion by categorizing conflicts with respect to their origin.
\begin{enumerate}
\item \emph{Ambiguity Conflicts} are conflicts that originate from an ambiguity
      in the underlying grammar. Such conflicts therefore indicate an actual
      problem with the grammar. We saw an example of such conflicts when we
      tried in Figure \ref{fig:Conflicts.ipynb}
      to describe the syntax of arithmetic expressions without the syntactic
      categories \textsl{product} and \textsl{factor}.

      In \simtextsc{Lark}, such conflicts are resolved by rewriting the grammar as shown in the previous
      section.
\item \emph{Look-Ahead Conflicts} are Reduce-Reduce conflicts where the grammar is
      technically unambiguous, but for which a look-ahead of one token is
      insufficient to resolve the conflict.
\item \emph{Mysterious Conflicts} arise only during the transition from LR states to \simtextsc{Lalr} states
      by merging states with the same core. These conflicts therefore occur
      precisely when the concept of an \simtextsc{Lalr} grammar is insufficient to describe the syntax of the
      language to be parsed.
\end{enumerate}
We will now look at the last two cases in detail and show ways in which the conflicts can be resolved.

\subsection{Look-Ahead Conflicts}
A Look-Ahead Conflict exists when the grammar is unambiguous, but a look-ahead of one
token is insufficient to decide which rule should be used for reduction. Figure
\ref{fig:lr-conflict.g} shows the grammar
\href{https://github.com/karlstroetmann/Formal-Languages/blob/master/Python/Chapter-10-Lark/03-Look-Ahead.ipynb}{\texttt{03-Look-Ahead.ipynb}}\footnote{
I found this grammar on the web on Pete Jinks' page at the address
\\[0.1cm]
\hspace*{1.3cm}
\href{http://www.cs.man.ac.uk/~pjj/cs212/ho/node19.html}{\texttt{http://www.cs.man.ac.uk/\symbol{126}pjj/cs212/ho/node19.html}}
\\[0.1cm]
.},
which, although unambiguous, does not possess the LR(1) property and is therefore certainly not an \simtextsc{Lalr}(1) grammar.

\begin{figure}[!ht]
\centering
\begin{minted}[ frame        = lines, 
                framesep     = 0.3cm, 
                bgcolor      = sepia,
                firstnumber  = 1,
                numbers      = left,
                numbersep    = -0.2cm,
                xleftmargin  = 0.8cm,
                xrightmargin = 0.8cm,
              ]{text}
    a : b "U" "V"
      | c "U" "W"
    b : "X"
    c : "X"
\end{minted}
\vspace*{-0.3cm}
\caption{An unambiguous grammar without the LR(1) property.}
\label{fig:lr-conflict.g}
\end{figure}

If we calculate the LR states of this grammar,
we find, among others, the following state:
\\[0.2cm]
\hspace*{1.3cm}
$\{ b \rightarrow \;\saquoted{X} \bullet: \saquoted{U},\; c \rightarrow \;\saquoted{X} \bullet: \saquoted{U} \}$.
\\[0.2cm]
Since the set of follow tokens is the same for both rules, we have a Reduce-Reduce conflict here.
This conflict is caused by the fact that the parser, with a look-ahead of only one token, cannot
decide whether an $\saquoted{X}$ should be interpreted as a $b$ or as a $c$, because this
is only decided when the character following $\saquoted{U}$ is read: If this is
a $\saquoted{V}$, then the rule
\\[0.2cm]
\hspace*{1.3cm}
$a \rightarrow b\; \saquoted{U}\; \saquoted{V}$
\\[0.2cm]
will be used overall, and consequently the $\saquoted{X}$ is to be interpreted as a $b$. If, however, the second token
after the $\saquoted{X}$ is a $\saquoted{W}$, then the rule to be used is
\\[0.2cm]
\hspace*{1.3cm}
$a \rightarrow c \;\saquoted{U}\; \saquoted{W}$
\\[0.2cm]
and consequently the $\saquoted{X}$ is to be read as $c$.

\simtextsc{Lark} refuses to build a parser for this grammar.  Instead, an exception of class
\texttt{GrammarError} is raised whose message is shown in Figure \ref{fig:Look-Ahead.error} on page
\pageref{fig:Look-Ahead.error}.  Note that the part of the message that shows the state is only printed if the
parser is created with \texttt{debug=True}.

\begin{figure}[!ht]
\centering
\begin{minted}[ frame        = lines, 
                framesep     = 0.3cm, 
                bgcolor      = sepia,
                firstnumber  = 1,
                numbers      = left,
                numbersep    = 0.4cm,
                xleftmargin  = 0.0cm,
                xrightmargin = 0.0cm,
              ]{text}
Reduce/Reduce collision in Terminal('U') between the following rules: 
    - <c : X>
    - <b : X>
    collision occurred in state: {
        <c : X * >
        <b : X * >
    }
\end{minted}
\vspace*{-0.3cm}
\caption{The error message produced for the grammar of Figure \ref{fig:lr-conflict.g}.}
\label{fig:Look-Ahead.error}
\end{figure}

\begin{figure}[!ht]
\centering
\begin{minted}[ frame        = lines, 
                framesep     = 0.3cm, 
                bgcolor      = sepia,
                firstnumber  = 1,
                numbers      = left,
                numbersep    = -0.2cm,
                xleftmargin  = 0.8cm,
                xrightmargin = 0.8cm,
              ]{text}
    a : b "V"
      | c "W"
    b : "X" "U"
    c : "X" "U"
\end{minted}
\vspace*{-0.3cm}
\caption{An LR(1) grammar equivalent to Figure \ref{fig:lr-conflict.g}.}
\label{fig:lr-conflict-resolved.g}
\end{figure}

The problem with this grammar is that it attempts to interpret an $\saquoted{X}$ either
as a $b$ or as a $c$ depending on the context. It is obvious how the problem can be
solved: If the context ``\texttt{U}'', which follows both $b$ and $c$, is included in
the rules for $b$ and $c$, then the conflict disappears, because then the
state in which the conflict previously occurred has the form
\\[0.2cm]
\hspace*{1.3cm}
$\{ b \rightarrow \;\saquoted{X}\; \saquoted{U}\bullet: \saquoted{V},\;
    c \rightarrow \;\saquoted{X}\; \saquoted{U}\bullet: \saquoted{W}
\}
$.
\\[0.2cm]
Here, the next token decides which rule we must reduce with in this state:
If the next token is a $\saquoted{V}$, we reduce with the rule
\\[0.2cm]
\hspace*{1.3cm}
$b \rightarrow \;\saquoted{X}\; \saquoted{U}$,
\\[0.2cm]
if, however, the next token is the letter $\saquoted{W}$, we take the rule
\\[0.2cm]
\hspace*{1.3cm}
$c \rightarrow \;\saquoted{X}\; \saquoted{U}\bullet: \saquoted{W}$ 
\\[0.2cm]
instead.  Figure
\ref{fig:lr-conflict-resolved.g} shows the correspondingly modified grammar, which you can find at
\\[0.2cm]
\hspace*{1.3cm}
\href{https://github.com/karlstroetmann/Formal-Languages/tree/master/Python/Chapter-10-Lark/04-Look-Ahead-Solved.ipynb}{\texttt{Python/Chapter-10-Lark/04-Look-Ahead-Solved.ipynb}}
\\[0.2cm]
on the web.  For this grammar, \simtextsc{Lark} builds a parser without complaining.  Applying the function
\texttt{dump\_states} to this parser shows the decisive state, which is reproduced in Figure
\ref{fig:Look-Ahead-Solved.state7} on page \pageref{fig:Look-Ahead-Solved.state7}.  The follow sets of the two
rules are now disjoint and therefore the action table of this state offers exactly one reduce action per
token: a \saquoted{V} reduces with the rule for $b$, while a \saquoted{W} reduces with the rule for $c$.

\begin{figure}[!ht]
\centering
\begin{minted}[ frame        = lines, 
                framesep     = 0.3cm, 
                bgcolor      = sepia,
                firstnumber  = 1,
                numbers      = left,
                numbersep    = -0.2cm,
                xleftmargin  = 0.8cm,
                xrightmargin = 0.8cm,
              ]{text}
    state 7
        <b : X U * >
        <c : X U * >

        V        reduce using rule <b : X U>
        W        reduce using rule <c : X U>
\end{minted}
\vspace*{-0.3cm}
\caption{The state that had caused the conflict, now free of conflicts.}
\label{fig:Look-Ahead-Solved.state7}
\end{figure}
\FloatBarrier

\remarkEng
There is a second way out: the grammar can be parsed with the \simtextsc{Earley} parser, which is the default
parser of \simtextsc{Lark}.  If the parser is created with
\\[0.2cm]
\hspace*{1.3cm}
\texttt{Lark(grammar, start='a', parser='earley', keep\_all\_tokens=True)}
\\[0.2cm]
then the grammar of Figure \ref{fig:lr-conflict.g} is accepted as it stands, because the
\simtextsc{Earley} algorithm does not need a bounded look-ahead.  The keyword argument
\texttt{keep\_all\_tokens} keeps the anonymous terminals in the parse tree, which would otherwise be filtered
out.  Notebook
\href{https://github.com/karlstroetmann/Formal-Languages/tree/master/Python/Chapter-10-Lark/03-Look-Ahead.ipynb}{\texttt{03-Look-Ahead.ipynb}}
demonstrates this: the \simtextsc{Earley} parser recognizes the 
\saquoted{X} of the string \aquoted{XUV} as a $b$ and the \saquoted{X} of the string \aquoted{XUW} as a $c$.
The price is a parser that is considerably slower than an \simtextsc{Lalr} parser.
\eox

\subsection{Mysterious Reduce-Reduce Conflicts}
A conflict is called a  \blue{mysterious reduce-reduce conflict} if the conflict results from the merger of
states that happens when we go from an LR parsing table to an \simtextsc{Lalr} parsing table.  The grammar in Figure
\ref{fig:Mysterious-Conflicts.ipynb} on page \pageref{fig:Mysterious-Conflicts.ipynb} is the same as the
grammar shown in Figure \ref{fig:lr-but-notlalr.g} on page \pageref{fig:lr-but-notlalr.g} in the previous
chapter.  Then we had seen that this grammar is an LR grammar, but not an \simtextsc{Lalr} grammar.  Let us see what happens
if we use \simtextsc{Lark} to generate the states for this grammar.  This example is discussed in the notebook
\\[0.2cm]
\hspace*{1.3cm}
\href{https://github.com/karlstroetmann/Formal-Languages/tree/master/Python/Chapter-10-Lark/05-Mysterious-Conflicts.ipynb}{\texttt{Python/Chapter-10-Lark/05-Mysterious-Conflicts.ipynb}}.

\begin{figure}[!ht]
\centering
\begin{minted}[ frame        = lines, 
                framesep     = 0.3cm, 
                bgcolor      = sepia,
                firstnumber  = 1,
                numbers      = left,
                numbersep    = -0.2cm,
                xleftmargin  = 0.8cm,
                xrightmargin = 0.8cm,
              ]{text}
    s : "v" a "y"
      | "w" b "y"
      | "v" b "z"
      | "w" a "z"

    a : X

    b : X

    X : "x"
\end{minted}
\vspace*{-0.3cm}
\caption{A grammar that generates a mysterious reduce-reduce conflict.}
\label{fig:Mysterious-Conflicts.ipynb}
\end{figure}
\vspace*{0.3cm}

When we compute the LR(1) states of this grammar, we find the two states
\\[0.2cm]
\hspace*{1.3cm}
$\bigl\{ a \rightarrow X \bullet: \saquoted{y},\;
         b \rightarrow X \bullet: \saquoted{z}
\bigr\}$
\quad and \quad
$\bigl\{ b \rightarrow X \bullet: \saquoted{y},\;
         a \rightarrow X \bullet: \saquoted{z}
\bigr\}$.
\\[0.2cm]
Taken by itself, neither of these two states has a conflict, since the follow sets of the respective rules are
disjoint.  However, these two states have the same core and are therefore merged when the
\simtextsc{Lalr} table is computed.  The resulting state has the form
\\[0.2cm]
\hspace*{1.3cm}
$\bigl\{ a \rightarrow X \bullet:\{\saquoted{y}, \saquoted{z} \},
         b \rightarrow X \bullet:\{\saquoted{y}, \saquoted{z} \}
\bigr\}$
\\[0.2cm]
and obviously has a reduce-reduce conflict if the next token is either \saquoted{y} or \saquoted{z}.

\simtextsc{Lark} performs this merger and therefore reports two reduce-reduce collisions, one for the token
\saquoted{y} and one for the token \saquoted{z}.  Figure \ref{fig:Mysterious-Conflicts.error} on page
\pageref{fig:Mysterious-Conflicts.error} shows the resulting error message.  In other words,
\simtextsc{Lark} really does construct \simtextsc{Lalr} tables and therefore rejects this grammar, as the
theory of the previous chapter predicts.

\begin{figure}[!ht]
\centering
\begin{minted}[ frame        = lines, 
                framesep     = 0.3cm, 
                bgcolor      = sepia,
                firstnumber  = 1,
                numbers      = left,
                numbersep    = 0.4cm,
                xleftmargin  = 0.0cm,
                xrightmargin = 0.0cm,
              ]{text}
Reduce/Reduce collision in Terminal('Z') between the following rules: 
    - <a : X>
    - <b : X>
    collision occurred in state: {
        <a : X * >
        <b : X * >
    }

Reduce/Reduce collision in Terminal('Y') between the following rules: 
    - <a : X>
    - <b : X>
    collision occurred in state: {
        <a : X * >
        <b : X * >
    }
\end{minted}
\vspace*{-0.3cm}
\caption{The error message for the grammar of Figure \ref{fig:Mysterious-Conflicts.ipynb}.}
\label{fig:Mysterious-Conflicts.error}
\end{figure}
\vspace*{0.3cm}

Note that the error message shows the merged state, but it does not show the follow sets of the two rules.
These have to be reconstructed by hand, as we have just done.

\remarkEng
A mysterious conflict cannot be repaired by rewriting the grammar in the way that we have used for the
look-ahead conflict of the previous subsection, for the grammar is already an LR(1) grammar: the defect is not
in the grammar but in the compression of the LR states into the \simtextsc{Lalr} states.  Hence, the only way
to parse this language with \simtextsc{Lark} is to use the \simtextsc{Earley} parser:
\\[0.2cm]
\hspace*{1.3cm}
\texttt{Lark(grammar, start='s', parser='earley', keep\_all\_tokens=True)}
\\[0.2cm]
The \simtextsc{Earley} algorithm builds no parse table at all and therefore never merges any states.  The
notebook shows that all four strings of the language are then parsed correctly: the \saquoted{x} in the middle
is recognized as an $a$ in \aquoted{vxy} and \aquoted{wxz}, but as a $b$ in \aquoted{vxz} and \aquoted{wxy}.
\eox

\exerciseEng
Use \simtextsc{Lark} to implement a \textsl{Python} parser that is able to evaluate formulas from propositional logic.
The language should support the operators
\begin{itemize}
\item ``\texttt{<->}'' (equivalence), 
\item ``\texttt{->}'' (implication),
\item ``\texttt{|}'' (disjunction), 
\item ``\texttt{\&}'' (conjunction), and 
\item ``\texttt{!}'' (negation).
\end{itemize}
The precedences and associativities are defined as follows:
\begin{enumerate}
\item The operator ``\texttt{!}''  binds stronger than the operator ``\texttt{\&}''.
\item The operator ``\texttt{\&}'' binds stronger than the operator ``\texttt{|}''.  
\item The operators ``\texttt{|}'' and ``\texttt{\&}'' are both left associative.  
\item The operator ``\texttt{->}'' associates to the right and binds weaker than the operator ``\texttt{|}''. 
\item The operator ``\texttt{<->}'' is not associative and binds weaker than the operator ``\texttt{->}''.
\end{enumerate}
Furthermore, the language should support parenthesis, variables, and the constants ``$\top$'' and ``$\bot$''.
Since \simtextsc{Lark} does not support precedence declarations, you have to introduce one syntactical variable
for every precedence level.  Create the parser with \texttt{parser='lalr'} and \texttt{strict=True} so that you
notice immediately if your grammar has a conflict.

\exerciseEng
Implement a parser (using \simtextsc{Lark}) that evaluates operations on sets of integers.
Your parser must accept strings containing integer sets and operators.
A set literal is defined as a comma-separated list of integers enclosed in curly braces.
An empty set is allowed.
\begin{itemize}
    \item Example: \texttt{\{1, 2, 3\}}
    \item Example: \texttt{\{\}}
\end{itemize}
You must support the following operators. The table is ordered from \textbf{highest binding power (top)} to \textbf{lowest binding power (bottom)}.

\begin{table}[h]
\centering
\begin{tabular}{lccl}

  \textbf{Precedence} & \textbf{Name} & \textbf{Math Symbol} & \textbf{Associativity} \\ 
4 & Power Set    &  $2^A$ & Right \\
3 & Intersection &  $A \cap B$ & Left \\
2 & Symmetric Difference & $A \,\Delta\, B$ & Left \\
1 & Union, Difference   & $A \cup B$, $A \setminus B$ & Left \\ 
\end{tabular}
\caption{Operator Precedence Table}
\label{tab:set_ops}
\end{table}

\noindent
Again, the precedences have to be encoded into the grammar.  Note that the parser has to be an
\simtextsc{Lalr} parser: if you use the default \simtextsc{Earley} parser, an ambiguous grammar will not be
rejected but will produce ambiguous parse trees instead.

\section{Check your Understanding}
\begin{enumerate}[(a)]
\item What do you have to do in order to see shift-reduce conflicts when using \simtextsc{Lark}?
\item How does \simtextsc{Lark} resolve a shift-reduce conflict?  How does it deal with a reduce-reduce
      conflict?
\item What is the effect of the keyword argument \texttt{strict=True}?
\item Are you able to inspect the \simtextsc{Lalr} states that \simtextsc{Lark} has computed?  Which keyword
      argument is needed for this?
\item Why is the string ``\texttt{1*2+3}'' parsed as ``\texttt{1*(2+3)}'' if the grammar does not
      specify the precedences of the arithmetical operators?
\item Given the precedences and associativities of a set of operators, can you write down a grammar that is an
      \simtextsc{Lalr}(1) grammar and that respects these precedences and associativities?
\item How does the grammar rule for an operator of a given precedence level depend on whether this operator is
      \blue{left associative}, \blue{right associative}, or \blue{non-associative}?
\item What is the purpose of the question mark in a rule of the form ``\texttt{?expr : \dots}'' and what is the
      purpose of an \blue{alias} that is introduced with the operator ``\texttt{->}''?

      Can these features introduce or remove a conflict?
\item Which three kinds of conflicts have we distinguished according to their origin?
\item How is the concept of a \blue{look-ahead conflict} defined?
\item How is the concept of a \blue{mysterious reduce-reduce conflict} defined?  
\end{enumerate}
    
%%% Local Variables: 
%%% mode: latex
%%% TeX-master: "formal-languages"
%%% End:
